% TOP_PROBLEM: {"id": 4700018, "title": "A piecewise-linear Hilbert sixteenth-type problem", "category_id": 11, "category": {"id": 11, "name": "dynamical_systems", "display_name": "Dynamical Systems", "description": "Problems about long-term behavior of deterministic systems, Hamiltonian dynamics, and periodic orbits.", "slug": "dynamical-systems", "order_index": 11, "created_at": "2026-07-31T15:26:25.671Z"}, "status": "partially_solved", "problem_number": "AMR-046-0018", "set_id": 13, "statusQualification": "The displayed inequalities are the 2020 source benchmarks; this entry asks the continuing extremal construction question, not whether those benchmarks are current records."}
% FIRST_POSTED: 2026-10-02
% ENTRY_KIND: statement_only
% UnsolvedMath ID 4700018; source catalogue code AMR-046-0018
% !TeX program = lualatex
% Definition and sources only; this file is not an LLM solution attempt.
\documentclass[11pt,a4paper]{article}
\usepackage[margin=25mm]{geometry}
\usepackage{fontspec}
\setmainfont{lmroman10-regular.otf}[BoldFont=lmroman10-bold.otf,
 ItalicFont=lmroman10-italic.otf,BoldItalicFont=lmroman10-bolditalic.otf]
\usepackage{amsmath,amssymb,mathtools}
\usepackage{unicode-math}
\setmathfont{latinmodern-math.otf}
\AtBeginDocument{\let\setminus\smallsetminus}
\usepackage{xurl,hyperref}
\hypersetup{colorlinks=true,urlcolor=blue}
\setcounter{secnumdepth}{0}
\setlength{\parindent}{0pt}
\setlength{\parskip}{5pt}
\setlength{\emergencystretch}{3em}
\begin{document}
\section*{A piecewise-linear Hilbert sixteenth-type problem}\label{4700018}
{\small UnsolvedMath ID 4700018 · AMR-046-0018 · Dynamical Systems · AMR Open Problem Lists}
\subsection{Definitions and mathematical statement}
Let $H(n)$ be the supremum of the number of isolated nonconstant periodic
orbits of a real planar polynomial vector field of total degree at most $n$.
For a second class, let a smooth embedded branch $\Sigma$ of a real
degree-$n$ algebraic curve separate the plane into two zones. Put an affine
linear vector field in each zone, allowing different fields on the two
sides. A crossing periodic orbit follows these fields and crosses $\Sigma$
transversely, with the two normal components having the same nonzero sign
at every crossing. No sliding segment or tangency is included.

Let $L(n)$ be the corresponding supremum of the number of isolated
crossing periodic orbits, over admissible separating branches and both
affine fields. The source asks to improve the lower bounds
\[
 H(2)\ge4,\quad H(3)\ge13,\quad H(n)\ge K n^2\log n,
 \qquad
 L(1)\ge3,\quad L(2)\ge4,\quad L(n)\ge\lfloor n/2\rfloor,
\]
where the asymptotic estimate uses an absolute $K>0$ for sufficiently large
$n$. These are the source's historical benchmarks, not assertions that
they remain the best published estimates.

The continuing construction problem is to determine how large these
quantities can be and improve available degree-dependent guarantees.
Each construction must verify isolation of every counted orbit and, for
$L(n)$, the prescribed two-zone geometry and transverse crossing condition.
Adding more zones or allowing a nonlinear vector field inside a zone
changes the class and does not establish a bound for $L(n)$.
\subsection{Short English statement}
Improve constructions with many limit cycles in polynomial systems and in two-zone piecewise-linear systems with algebraic separation.
\subsection{Sources}
\begin{itemize}
\item[S1] ULAM.AI and the UnsolvedMath Contributors, supplied problems.json, record 4700018 (AMR-046-0018), AMR Open Problem Lists. Catalogue identity and source transcription; CC BY 4.0.
\url{https://huggingface.co/datasets/ulamai/UnsolvedMath}.
\item[S2] Gasull - Some open problems in low dimensional dynamical systems (2020), Problem environment 18: A piecewise-linear Hilbert sixteenth-type problem. Original source indexed by AMR-046-0018.
\url{https://arxiv.org/abs/2012.02524}.
\end{itemize}
\end{document}
